%% kapa: exam=1 round=26 year=2015 period=1 form=A subject=civil_law subject_ko=민법 seq=20 number=20
%% source: Q-Net 감정평가사 기출문제 (cst00309) artlSeq=5138001 file=extracted/1008/5138001/제26회_감정평가사 1차_기출/제26회_감정평가사 1차_기출/1교시_감평_A형.pdf
%% conversion: kordoc 4.24.51 -> tools/gen1.py (auto); review=proofread-sonnet; flags=none
\begin{question}{20}
소멸시효 중단에 관한 설명으로 옳은 것을 모두 고른 것은? (다툼이 있으면 판례에 따름)
\begin{framedtext}
ㄱ.채무자가 제기한 채무부존재확인소송에서 채권자가 피고로서 응소하여 적극적으로 권리를 주장하고 그것이 법원에 의해 받아들여진 경우, 채권의 소멸시효가 중단된다. \par
ㄴ.비법인사단의 대표자가 총회결의에 따라 총유물을 매도하여 소유권이전등기를 해주기 위해 매수인과 함께 법무사 사무실을 방문한 행위는, 소유권이전등기청구권의 소멸시효 중단의 효력이 있는 승인에 해당한다. \par
ㄷ.재판상 청구로 인하여 중단된 시효는 재판이 시작된 때부터 새로 진행된다.
\end{framedtext}
\begin{choices}
  \item ㄱ
  \item ㄴ
  \item ㄱ, ㄴ
  \item ㄱ, ㄷ
  \item ㄴ, ㄷ
\end{choices}
\end{question}
\answer{3}
